%source mathadata
\begin{multicols}{2}
			\begin{tikzpicture}[scale=0.75]
				\tikzset{
					grid major/.style = {very thin, gray},
					grid minor/.style = {very thin, lightgray}
				}
				\draw[grid major] (-1, -1) grid[step=1] (8, 7);
				\draw[->] (-1,0) -- (8.5,0) node[below] {x};
				\draw[->] (0,-1) -- (0,7.5) node[right] {y};
				\foreach \x in {0,1,...,8} \draw (\x,0.2) -- (\x,-0.2) node[below] {\x};
				\foreach \y in {0,1,...,7} \draw (0.2,\y) -- (-0.2,\y) node[left] {\y};

				% Points de l'ensemble E (points blancs)
				\node at (3,5) [circle,fill,inner sep=3pt,label=above:{$E_1$}] {};
				\node at (5,4) [circle,fill,inner sep=3pt,label=above right:{$E_2$}] {};

				% Points de l'ensemble F (points noirs)
				\node at (5,1) [circle,fill=white,thick,draw,inner sep=3pt,label=below right:{$F_1$}] {};
				\node at (3,2) [circle,fill=white,thick,draw,inner sep=3pt,label=below:{$F_2$}] {};

				% Points A, B, C
				\node at (4,2) [cross out, draw, red, inner sep=2pt, line width=0.8pt, label=below right:{$A$}] {};
				\node at (2,3) [cross out, draw, red, inner sep=2pt, line width=0.8pt, label=above:{$B$}] {};
				\node at (0,1) [cross out, draw, red, inner sep=2pt, line width=0.8pt, label=left:{$C$}] {};

				% Droite D
				\draw[domain=-1:8, smooth, variable=\x, blue, thick] plot ({\x}, {0.5*\x+3});
				\node[font=\small, blue] at ({8.5}, {0.5*8.5+3}) {$D$};

				% Droites D1, D2, D3
				\draw[domain=-1:8, smooth, variable=\x, green!70!black, dashed] plot ({\x}, {0.5*\x});
				\node[font=\small, green!70!black] at ({8.5}, {0.5*8.5}) {$D_1$};

				\draw[domain=-1:8, smooth, variable=\x, orange, dashed] plot ({\x}, {0.5*\x+2});
				\node[font=\small, orange] at ({8.5}, {0.5*8.5+2}) {$D_2$};

				\draw[domain=-1:8, smooth, variable=\x, purple, dashed] plot ({\x}, {0.5*\x+1});
				\node[font=\small, purple] at ({8.5}, {0.5*8.5+1}) {$D_3$};

			\end{tikzpicture}

			\begin{enumerate}
				\item Voir le graphique ci-contre.

				\item Les trois droites sont parallèles à $D$, donc de la forme $y = 0{,}5x + p$.
				      \begin{itemize}
					      \item $D_1$ passe par $A(4;2)$ : $2 = 0{,}5 \times 4 + p$ donc $p = 0$.

					            Équation : $\boldsymbol{y = 0{,}5x}$

					      \item $D_2$ passe par $B(2;3)$ : $3 = 0{,}5 \times 2 + p$ donc $p = 2$.

					            Équation : $\boldsymbol{y = 0{,}5x + 2}$

					      \item $D_3$ passe par $C(0;1)$ : $1 = 0{,}5 \times 0 + p$ donc $p = 1$.

					            Équation : $\boldsymbol{y = 0{,}5x + 1}$
				      \end{itemize}

				\item Par comptage graphique :
				      \begin{itemize}
					      \item $D_1$ : 1 erreur (le point noir $(3;2)$ est au-dessus)
					      \item $D_2$ : 1 erreur (le point blanc $(5;4)$ est en dessous)
					      \item $D_3$ : 0 erreur
				      \end{itemize}

				      \textbf{Réponse : $D_3$ sépare les deux paires de points.}

				\item Vérification par le calcul pour $D_3 : y = 0{,}5x + 1$ :
				      \begin{itemize}
					      \item $E_1(3;5)$ : $0{,}5 \times 3 + 1 = 2{,}5 < 5$ \\  $E_1$ est au-dessus de $D_3$
					      \item $E_2(5;4)$ : $0{,}5 \times 5 + 1 = 3{,}5 < 4$ \\  $E_2$ est au-dessus de $D_3$
					      \item $F_1(5;1)$ : $0{,}5 \times 5 + 1 = 3{,}5 > 1$ \\  $F_1$ est en dessous de $D_3$
					      \item $F_2(3;2)$ : $0{,}5 \times 3 + 1 = 2{,}5 > 2$ \\  $F_2$ est en dessous de $D_3$
				      \end{itemize}

			\end{enumerate}
		\end{multicols}
